\chapter{Conclusion}\label{ch:conclusion}

This thesis demonstrated a technical feasibility pipeline for a low-cost sacrum-mounted wearable that detects pelvis-region movement during prolonged sitting and delivers a local haptic cue when such movement is absent. The contribution is not a validated behavioural intervention, but an integrated sensing-to-cue prototype that connects IMU sensing, labelled mobile data collection, machine-learning detection, compact embedded deployment, and vibrotactile feedback in one testable system. The work was organised around four linked objectives: prototype implementation, static-sitting versus pelvic-rotation detection, Shimmer/XIAO trend-level sensing agreement, and compact embedded deployment with haptic reminder support.

For Objective~1, the thesis shows that a compact wearable research prototype can integrate the required sensing-to-cue pipeline. The final system combines a XIAO nRF54L15 Sense sensor node worn near the sacrum, the embedded LSM6DS3TR-C IMU, BLE streaming, Android labelled field recording, PC researcher-side monitoring, compact on-device classification, a DRV2605L haptic driver, and a coin vibration motor. This demonstrates the prototype as an integrated research platform rather than only a sensor logger or an offline analysis workflow. The main remaining engineering weaknesses are mechanical robustness, battery-contact reliability, and the absence of measured battery lifetime.

For Objective~2, labelled pelvis-region IMU data could distinguish static sitting from instructed pelvic rotation in the preliminary pilot dataset. In the person-independent split using participants P1--P4 for training and participant P5 for final testing, the \(k=7\) KNN classifier achieved an F1-score of 0.957 for pelvic rotation, detecting 155 of 165 pelvic-rotation windows and producing 4 false positives among 3{,}497 static-sitting windows. This supports detectability of the instructed movement pattern in the collected feature space. It does not establish general-purpose sitting-behaviour recognition, because the evaluation was window-level, based on a small instructed dataset, and tested on one held-out participant.

For Objective~3, the custom XIAO-based wearable captured motion trends that were sufficiently consistent with a Shimmer reference for the binary movement-detection task studied here. In the co-located comparison, Pearson correlations across acceleration and gyroscope axes ranged from 0.816 to 0.923, with estimated lag between \(-0.020\) and 0~s. In the two valid pilot field-worn paired recordings, maximum cross-correlations reached 0.915--0.981 after lag correction. These results support trend-level use of the XIAO signal for this thesis, but they do not establish clinical pelvic-angle measurement or full biomechanical validation. The limited paired Shimmer coverage and the 1.2--1.3~s field lag remain important evidence boundaries.

For Objective~4, the selected detector could be compressed and deployed on the target microcontroller. The compact 12-feature, 128-prototype KNN achieved F1 = 0.959 on the same held-out participant and ran in 2.73~ms per 2~s window on the XIAO nRF54L15 Sense, using 56.3~KB of flash and 9.3~KB of RAM. The compact model demonstrates deployability for this instructed-movement task, not model optimality. Timing feasibility also does not yet prove battery-life feasibility or behavioural effectiveness.

The thesis therefore demonstrates the full technical chain required for a local wearable reminder: pelvis-region sensing, labelled mobile data collection, windowed feature extraction, classifier comparison, reference-sensor trend comparison, prototype compression, embedded inference, inactivity timing, and haptic actuation. This is broader than an offline classification experiment, but narrower than a behavioural intervention trial or a finished engineering product.

This boundary is essential. The prototype can detect the instructed movement in the collected pilot data and trigger a local tactile cue without continuous smartphone or PC-based processing. It does not yet show that such cues change long-term sitting behaviour, improve posture, reduce health risk, or remain reliable during unsupervised daily wear. Future work should therefore collect more natural seated movement from a larger and more diverse cohort, improve mechanical and battery reliability, measure power consumption and battery life, and run a dedicated user study on haptic reminder awareness, acceptability, and behavioural response.

Overall, the thesis turns the idea of interrupting prolonged static sitting with local wearable feedback into a measurable prototype pipeline. Its main contribution is the integration and preliminary validation of the sensing, detection, deployment, and cue-delivery chain. Its main boundary is equally clear: the system is technically feasible for the instructed pilot task, but behavioural effectiveness, clinical validity, battery-life feasibility, and unsupervised daily-wear robustness remain to be demonstrated.
